\section{排版核对清单示例}
\label{app:checklist}
以下清单只展示跨页表格的用法。条目是通用检查提示，不代表学校完整验收表，也不表示本项目已经通过所有正式格式检查。长表不使用浮动表格环境；跨页后自动重复表头。

{\zihao{5}
\begin{longtable}{p{0.07\textwidth}p{0.22\textwidth}p{0.58\textwidth}}
\caption{写作与排版检查提示}\label{tab:checklist}\\
\toprule
序号 & 检查对象 & 检查提示 \\
\midrule
\endfirsthead
\multicolumn{3}{c}{表\thetable\quad 写作与排版检查提示（续）}\\
\toprule
序号 & 检查对象 & 检查提示 \\
\midrule
\endhead
\midrule
\multicolumn{3}{r}{续下页}\\
\endfoot
\bottomrule
\endlastfoot
1 & 规范版本 & 核实当届文件是否继续采用所附二〇二四届手册的规则。\\[4pt]
2 & 学院要求 & 记录已确认适用的补充要求，不以学院名称推断版式。\\[4pt]
3 & 打印封面 & 使用经确认的学校底稿，检查是否需要另附打印扉页。\\[4pt]
4 & 基本信息 & 替换示例姓名、学号、学院、专业、导师和日期。\\[4pt]
5 & 中英文题目 & 检查两种题目的含义对应关系与长题目的换行。\\[4pt]
6 & 中文摘要 & 核对摘要内容及字数，避免写入正文尚未支持的结论。\\[4pt]
7 & 英文摘要 & 核对译文、术语、标点以及关键词与中文的对应关系。\\[4pt]
8 & 关键词 & 选择反映核心内容的词语，并检查数量和分隔方式。\\[4pt]
9 & 自动目录 & 编译至交叉引用稳定后，核对标题和目录页码。\\[4pt]
10 & 标题层级 & 使用三级以内的语义层级，不手写编号和手动字号。\\[4pt]
11 & 正文字体 & 在正式字体方案下检查宋体、黑体和西文字体。\\[4pt]
12 & 页面边距 & 测量输出页面，并核对装订线是否与左边距叠加。\\[4pt]
13 & 页眉页脚 & 检查校名字标的尺寸、位置和右下角页码。\\[4pt]
14 & 数学符号 & 在首次使用时解释符号、单位和适用范围。\\[4pt]
15 & 公式编号 & 检查跨节编号方式和正文对公式的自动引用。\\[4pt]
16 & 图片内容 & 检查清晰度、图例、坐标轴名称和图片来源。\\[4pt]
17 & 表格内容 & 核对表头、计量单位、有效数字和数据来源。\\[4pt]
18 & 图表题名 & 核对图题在下、表题在上，以及题名与正文的对应。\\[4pt]
19 & 文献脚注 & 核实引用键、书目信息和引用的实际页码。\\[4pt]
20 & 网络来源 & 区分发布日期与访问日期，保留可核查的网址。\\[4pt]
21 & 文献排序 & 核对中文组在前，以及中文作者的显式拼音排序键。\\[4pt]
22 & 跨页对象 & 查看长表的续表头、页脚提示和最后一页的表底线。\\[4pt]
23 & 附录材料 & 检查补充材料是否必要、可追溯并在正文得到说明。\\[4pt]
24 & 最终文件 & 清理示例内容后，重新编译并逐页检查最终 PDF。\\[4pt]
25 & 术语一致性 & 对同一概念使用一致译名，在首次出现时解释缩写。\\[4pt]
26 & 数值一致性 & 对照正文、表格和公式，检查同一数据是否前后一致。\\[4pt]
27 & 图文对应 & 确认每幅图和每张表都在正文中得到讨论。\\[4pt]
28 & 研究结论 & 检查结论是否由已经呈现的证据支持，避免扩大适用范围。\\[4pt]
29 & 引用版本 & 比较书目版本与实际阅读版本，核对具体引用页码。\\[4pt]
30 & 修改记录 & 保存重要版式变更的依据和日期，便于后续追溯。\\[4pt]
31 & 示例清理 & 搜索示例作者、占位学号及人工构造标记，逐项替换。\\[4pt]
32 & 文件归档 & 保存最终源文件、参考文献库、所需图片和输出 PDF。\\[4pt]
\end{longtable}
}
